\chapter{Lung Lobectomy}
\section*{Overview}
A lobectomy removes one lobe of the lung and is a major form of anatomic lung resection.
\section*{Why It Is Done}
It is commonly performed for localized lung cancer and may be used for selected infections, bronchiectasis, or other destructive lung disease.
\section*{How It Is Performed}
The lobe and associated vessels and bronchus are divided through an open thoracotomy or minimally invasive VATS or robotic approach. Lymph nodes are often sampled.
\section*{Risks and Recovery}
Air leak, bleeding, infection, pneumonia, respiratory failure, arrhythmia, blood clots, and reduced lung function are possible. Chest-tube care, breathing exercises, mobilization, and pathology review are required.
\section*{References}
\begin{enumerate}\item National Cancer Institute. Lung Surgery. 2025.\item Current Medical Diagnosis and Treatment. McGraw-Hill; 2025.\end{enumerate}
\clearpage
